\input{configTD}
\usepackage{listings}

\lstset{
    basicstyle=\ttfamily\small,
    breaklines=true,
    frame=single,
    columns=fullflexible,
    keepspaces=true
}

\title{Analyse Numérique TP1 - Euler, visualisation et pendule}

\author{}
\date{}

\newenvironment{solution}
    {\par\vspace{0.5em}\begin{mdframed}[linewidth=0.5pt]\par}
    {\end{mdframed}\par\vspace{0.5em}}

\begin{document}
\sffamily
\maketitle

\section*{Objectifs}
\begin{itemize}
    \item Implémenter la méthode d'Euler sur une EDO simple.
    \item Comparer solution exacte et solution approchée.
    \item Visualiser l'évolution temporelle dans un notebook Jupyter.
    \item Modifier un code existant pour ajouter un effet physique simple : le frottement.
\end{itemize}

\section*{Partie 1 - Une EDO simple}
On considère le problème
\[
\begin{cases}
y'(t)=-ky(t),\\
y(0)=y_0,
\end{cases}
\qquad k>0.
\]
La solution exacte est
\[
y(t)=y_0e^{-kt}.
\]

\begin{enumerate}
    \item Montrer que le schéma d'Euler explicite s'écrit
    \[
    y_{n+1}=y_n-hky_n.
    \]

    \item Compléter et exécuter le code suivant dans un notebook.
\end{enumerate}

\begin{lstlisting}[language=Python]
import numpy as np
import matplotlib.pyplot as plt

def euler_explicit(f, t0, y0, h, n):
    t = np.zeros(n + 1)
    y = np.zeros(n + 1)
    t[0] = t0
    y[0] = y0

    for i in range(n):
        t[i + 1] = t[i] + h
        y[i + 1] = y[i] + h * f(t[i], y[i])

    return t, y

k = 0.8
y0 = 5.0
h = 0.25
n = 20

f = lambda t, y: -k * y
t, y_num = euler_explicit(f, 0.0, y0, h, n)
y_exact = y0 * np.exp(-k * t)

plt.figure(figsize=(8, 4))
plt.plot(t, y_exact, label="solution exacte")
plt.plot(t, y_num, "o-", label="Euler")
plt.xlabel("t")
plt.ylabel("y")
plt.grid(True)
plt.legend()
plt.show()
\end{lstlisting}

\begin{enumerate}[resume]
    \item Tester plusieurs pas $h$ et commenter l'évolution de l'erreur.

    \item Ajouter dans une nouvelle cellule le calcul de l'erreur maximale
    \[
    \max_i |y(t_i)-y_i|.
    \]
\end{enumerate}

\section*{Partie 2 - Première animation Jupyter}
L'objectif est d'animer le point numérique se déplaçant au cours du temps sur le graphe précédent.

\begin{enumerate}
    \item Importer \texttt{FuncAnimation} depuis \texttt{matplotlib.animation}.

    \item Construire une animation qui fait apparaître progressivement les points $(t_i,y_i)$.

    \item Exporter l'animation au format GIF ou l'afficher dans le notebook.
\end{enumerate}

\noindent On pourra partir du squelette suivant.

\begin{lstlisting}[language=Python]
from matplotlib.animation import FuncAnimation

fig, ax = plt.subplots(figsize=(8, 4))
ax.plot(t, y_exact, label="solution exacte")
line_num, = ax.plot([], [], "o-", label="Euler")
ax.set_xlim(t[0], t[-1])
ax.set_ylim(0, 1.1 * y0)
ax.grid(True)
ax.legend()

def update(frame):
    line_num.set_data(t[:frame + 1], y_num[:frame + 1])
    return (line_num,)

anim = FuncAnimation(fig, update, frames=len(t), interval=200, blit=True)
plt.show()
\end{lstlisting}

\section*{Partie 3 - Pendule sans frottement}
On considère maintenant le système du pendule sous forme simplifiée
\[
\theta''(t)+\omega_0^2\sin(\theta(t))=0.
\]
En posant
\[
v(t)=\theta'(t),
\]
on obtient le système du premier ordre
\[
\begin{cases}
\theta'(t)=v(t),\\
v'(t)=-\omega_0^2\sin(\theta(t)).
\end{cases}
\]

\noindent Exécuter puis commenter le code suivant.

\begin{lstlisting}[language=Python]
import numpy as np
import matplotlib.pyplot as plt

omega0 = 1.5
h = 0.02
n = 1500

t = np.zeros(n + 1)
theta = np.zeros(n + 1)
v = np.zeros(n + 1)

theta[0] = 0.8
v[0] = 0.0

for i in range(n):
    t[i + 1] = t[i] + h
    theta[i + 1] = theta[i] + h * v[i]
    v[i + 1] = v[i] - h * omega0**2 * np.sin(theta[i])

plt.figure(figsize=(8, 4))
plt.plot(t, theta)
plt.xlabel("t")
plt.ylabel("theta(t)")
plt.grid(True)
plt.show()

plt.figure(figsize=(5, 5))
plt.plot(theta, v)
plt.xlabel("theta")
plt.ylabel("v")
plt.grid(True)
plt.show()
\end{lstlisting}

\section*{Partie 4 - A vous de jouer : ajouter des frottements}
On remplace maintenant le modèle par
\[
\theta''(t)+c\,\theta'(t)+\omega_0^2\sin(\theta(t))=0,
\qquad c>0.
\]

\begin{enumerate}
    \item Écrire le nouveau système du premier ordre.

    \item Modifier le code précédent pour prendre en compte le coefficient de frottement $c$.

    \item Tester plusieurs valeurs de $c$.

    \item Comparer l'allure de $\theta(t)$ et du portrait de phase avec et sans frottement.
\end{enumerate}

\section*{Partie 5 - Bonus}
\begin{enumerate}
    \item Ajouter une animation du pendule dans le plan.

    \item Remplacer Euler explicite par une méthode de Runge-Kutta d'ordre 2 ou 4.

    \item Proposer une interprétation en lien avec une vibration amortie d'un élément de structure.
\end{enumerate}

\section*{Compte rendu attendu}
\begin{itemize}
    \item Le notebook exécute sans erreur.
    \item Les graphes sont commentés.
    \item L'effet du pas de temps est discuté.
    \item La modification avec frottement est présentée et interprétée.
\end{itemize}

\end{document}
